\documentclass[10pt]{scrartcl}
\usepackage{amsmath,amsthm,amssymb,mathtools,xfrac}
\usepackage{enumerate}
\usepackage[utf8]{inputenc}
%\usepackage{german}
\usepackage{a4wide}
\usepackage{graphics}
\usepackage{graphicx}
\usepackage[T1]{fontenc}
\usepackage{epsfig}
\usepackage{libertine}
\usepackage[pdftex]{hyperref}
\usepackage{refcount}
\usepackage{framed}
\usepackage{xcolor}
\usepackage{multicol}
\usepackage{enumitem}
\usepackage{caption}
\usepackage{gensymb}
\usepackage[export]{adjustbox}

\topmargin-.5in \textheight9in \oddsidemargin0in \textwidth6.5in

\begin{document}


\begin{minipage}{.5\textwidth}
{\scshape Department of Mathematics and Informatics \\
University of Cologne\\}
\small Prof. Dr. Gregor Gassner \\
\end{minipage}
\begin{minipage}{.3\textwidth}
{\scshape Department of Mathematics \\
Link\"{o}ping University\\}
\small Andrew R. Winters \\
\end{minipage}
\hfill
\begin{minipage}{.15\textwidth}
\begin{flushright}
%\includegraphics[width=3cm]{figs/logo-uni-koeln.jpg}\\
\today
\end{flushright}
\end{minipage}

\begin{center}
{\Large\bfseries
An Introduction to Climate Modeling}\\
\vspace{0.2cm}
{\bfseries Milestone 3 - Solving a zero-dimensional Energy Balance Model}
\end{center}


\section{Solving a zero-dimensional Energy Balance Model}


In the previous milestones you have already implemented some routines that are part of an energy balance model. Now you have the opportunity to combine your results into a zero-dimensional ($0D$) energy balance model (EBM) and solve the model using an ODE solver. Note that for zero-dimensional EBMs, spatial dependence is neglected, and $T$ is a global time-dependent averaged surface temperature. Let the following equation be given for the energy balance model:

\begin{align}
\overline{C}\frac{\partial {T}_A(t)}{\partial t}+A+B{T}_A(t)=\overline{S_\text{sol}}(t),
\label{EBM}
\end{align}
where $\overline{C}$ is the mean heat capacity, ${T}_A$ is an approximation to the area-average of the time dependent temperature, $A$ the $CO_2$ dependent radiative cooling, $B = 2.15$ the feedback effects such as water vapor cycles, lapse rate and cloud cover, $\overline{S_\text{sol}}(t)$ is the area-averaged solar forcing.
To solve the zero-dimensional EBM you should proceed as described in the next steps.

\begin{enumerate}
\item Implement a function \small\texttt{\detokenize{calc_area}}\normalsize, which determines the area fraction of each grid cell as a function of latitude, with an equidistant step size $h=\frac{\pi}{64}$ for a $128 \times 65$ grid, and returns the values as a vector.
\item Implement a \small\texttt{\detokenize{calc_mean}} \normalsize function, using the results from the first task, which takes a matrix of values of a given parameter like albedo or solar forcing for each grid cell and returns the mean value over the whole area.
\item Write a function \small\texttt{\detokenize{calc_radiative_cooling_co2}}\normalsize, which returns the $CO_2$ dependent radiative cooling~$A$ for a given $CO_2$ concentration $c$. As default value use the $CO_2$ concentration of the year 1950: $c_0 = 315.0$. The radiative cooling $A$ is given by
\begin{align}
A (c) = 210.3-5.35 \log\left(\frac{c}{c_0}\right).
\label{A}
\end{align}
\item In a function \small\texttt{\detokenize{compute_equilibrium}}\normalsize, solve the above equation (\ref{EBM}) using the forward Euler method and calculate an average temperature per year until equilibrium is reached.
\item Write a new function or adapt your existing one to add an option to solve the equation with the backward Euler method instead. Plot your results.

\end{enumerate}



 \section{Bonus}

 \begin{enumerate}
 \item Calculate the temperature evolution with a function \small\texttt{\detokenize{calc_evolution}} \normalsize similar to the function in task 4. This time for the parameter $A$ use transient $CO_2$ concentrations that you can take from the input file \small\texttt{\detokenize{co2_nasa.dat}} \normalsize containing NASA $CO_2$-Data from 1958 to 2025 with measurements taken from the Mauna Loa Observatory in Hawaii\footnote{https://gml.noaa.gov/ccgg/trends/mlo.html}.
Each row of this $CO_2$ data lists:\\[0.2cm]
year, month, decimal date, monthly average, seasonal average, \# measurements, std. deviation, uncertainty\\[0.2cm]
For the last three entries values of $-1, -9.99, -0.99$ mean ``missing data" and should be ignored.

\newpage

 This means that you need to solve the following equation with the backward Euler method:

 \begin{align}
 \overline{C}\frac{\partial T_A(t)}{\partial t} + A(t) + BT_A(t) = \overline{S_\text{sol}}(t),
 \label{EBM2}
 \end{align}
 with
 \begin{align}
A(t) = 210.3-5.35 \log\left(\frac{c}{\bar{c}_0(t)}\right),
\label{A1}
\end{align}
where $\bar{c}_0(t)$ is the average $CO_2$ concentration for a given year.

\item Use this transient $CO_2$ data to compute the average temperature for the year 2020.

\item Create a function \small\texttt{\detokenize{calc_transient_0DEBM}} \normalsize that computes and saves the average temperature and annual temperatures for each year for which the transient $CO_2$ data is available.

 \end{enumerate}

\pagebreak

\section{Control Solutions}


\begin{figure}[!htb]
\includegraphics[width=0.7\linewidth,center]{../../website/_assets/milestone_results/milestone3_forward.png}
\caption{EBM results calculated with forward Euler method (results calculated with backward Euler method are similar)}
\end{figure}



\end{document}
